\section{A certified unique saddle for every proportional reflected moment}
\label{gs:sec:main}

Put
\[
 f(x)=\log\Gamma(x),\qquad F(x)=-\psi(x),\qquad P(x)=\psi_1(x)
 \qquad(0<x<1).
\]
All three functions are strictly positive. The local log-gamma series
used below is classical; the new assertion here is the global
monotonicity inequality, with its explicit exact certificate, and its
application to the proportional moment regime. The monotonicity proof
is \emph{computer assisted}: elementary analysis controls both endpoints,
and a finite rational interval certificate controls the compact middle.
It uses no numerical sampling, floating-point special functions, or
unproved decimal values as a proof step. We prove the reflected-ratio
inequality directly; no priority claim is made for this inequality or
its symmetric special case.

Ordinary convexity of $f$ does not furnish a shortcut. For example,
the positive decreasing convex function $f_*(x)=x^{-2}-1$ has the
analogous ratio $J_*(x)=(2-x)/(1+x)$, which decreases strictly.
Moreover $\log\log\Gamma(x)$ is convex sufficiently near zero, so
global logarithmic concavity is unavailable here. The argument must
use more of the structure of the gamma function.

\begin{theorem}[Certified global log-gamma slope ratio]
\label{gs:thm:monotonicity}
Define
\[
 J(x)=\frac{F(x)f(1-x)}{F(1-x)f(x)}.
\]
Then $J$ is a strictly increasing real-analytic bijection from $(0,1)$
onto $(0,\infty)$. More quantitatively,
\begin{equation}\label{gs:eq:global-derivative}
 \frac{J'(x)}{J(x)}>\frac{176}{225}\qquad(0<x<1).
\end{equation}
Also $J(1-x)=J(x)^{-1}$, $J(1/2)=1$, and
$J(x)\sim1/\log(1/x)$ as $x\downarrow0$.
\end{theorem}

\subsection{Elementary endpoint proof}

Differentiating the displayed definition gives
\begin{equation}\label{gs:eq:H}
 \mathcal H(x):=\frac{J'(x)}{J(x)}
 =\frac{F(x)}{f(x)}+\frac{F(1-x)}{f(1-x)}
  -\frac{P(x)}{F(x)}-\frac{P(1-x)}{F(1-x)}.
\end{equation}
In particular $\mathcal H(1-x)=\mathcal H(x)$.
The familiar convergent expansions \cite[\S5.7]{joint:DLMF}, valid
for $|x|<1$ with the indicated real logarithm for $x>0$, are
\begin{align}
 f(x)&=-\log x-\gamma x+
               \sum_{n\ge2}\frac{(-1)^n\zeta(n)}n x^n,\notag\\
 f(1-x)&=\gamma x+\sum_{n\ge2}\frac{\zeta(n)}n x^n,\notag\\
 F(x)&=\frac1x+\gamma+
               \sum_{n\ge2}(-1)^{n-1}\zeta(n)x^{n-1},\notag\\
 F(1-x)&=\gamma+\sum_{n\ge2}\zeta(n)x^{n-1},\notag\\
 P(x)&=\frac1{x^2}+\sum_{n\ge2}(-1)^n(n-1)\zeta(n)x^{n-2},\notag\\
 P(1-x)&=\sum_{n\ge2}(n-1)\zeta(n)x^{n-2}.
 \label{gs:eq:series}
\end{align}
The differentiated series converge uniformly on every smaller closed
disk. Only the elementary bounds $1/2<\gamma<3/5$ and
$\zeta(n)\le\zeta(2)<2$ are needed in the endpoint argument; the
certificate below also independently encloses these constants.

For $0<x\le1/16$, rewrite \eqref{gs:eq:H} as
\[
 \mathcal H(x)=\frac{F(x)}{f(x)}-\frac{P(1-x)}{F(1-x)}
 +\left(\frac{F(1-x)}{f(1-x)}-\frac1x\right)
 +\left(\frac1x-\frac{P(x)}{F(x)}\right).
\]
Both parenthesized terms are nonnegative. For the first,
\[
 xF(1-x)-f(1-x)
 =\sum_{n\ge2}\frac{n-1}{n}\zeta(n)x^n>0.
\]
For the second, the convergent series yields
\begin{align*}
 F(x)-xP(x)
 &=\gamma+\sum_{k\ge1}(-1)^k(k+1)\zeta(k+1)x^k\\
 &>\frac12-2\bigl((1-x)^{-2}-1\bigr)
 \ge\frac12-\frac{62}{225}=\frac{101}{450}>0.
\end{align*}
Also
\[
 F(x)>\frac1x+\frac12-\frac{2x}{1-x}>\frac1x.
\]
Log-convexity of $\Gamma$ between one and two gives
$0<f(x)\le-\log x$. Since $x\log(1/x)$ increases on $(0,1/16]$
and $\log16<3$, it follows that
\[
 \frac{F(x)}{f(x)}>\frac1{x\log(1/x)}>\frac{16}{3}.
\]
The inequality $\log16<3$ follows, for example, from the strict
trapezoidal upper bound $\log2<3/4$ for the integral of $1/t$ on
$[1,2]$. Finally,
\[
 \frac{P(1-x)}{F(1-x)}
 <\frac{4}{(1-x)^2}\le\frac{1024}{225}.
\]
Combining the four bounds proves
$\mathcal H(x)>16/3-1024/225=176/225$ on $(0,1/16]$.
Reflection proves the same bound on $[15/16,1)$.

\subsection{The compact interval certificate}

\begin{lemma}[Exact finite compact certificate]\label{gs:lem:certificate}
For every $x\in[1/16,1/2]$,
\[
 \mathcal H(x)>\frac74.
\]
\end{lemma}

\begin{proof}[Proof by a finite rational certificate]
The dependency-free program \texttt{certify\_gamma\_saddle.py}
recomputes the following enclosure procedure from scratch.
Every interval endpoint is an integer multiple of $2^{-100}$.
Addition is exact; multiplication, inversion and division round
outward using integer floor and ceiling. Thus all intermediate
intervals contain the corresponding exact real quantities, by induction
on the arithmetic operations. No stored interval endpoint needs to be
trusted by the verifier.

For $2\le n\le40$, enclose $\zeta(n)$ using the first $64$ terms
and the integral-test bounds
\[
 \sum_{k=1}^{64}k^{-n}+
       \frac{65^{1-n}}{n-1}
 \le\zeta(n)\le
 \sum_{k=1}^{64}k^{-n}+
       \frac{64^{1-n}}{n-1}.
\]
For $\gamma$, use
\[
 \gamma=\sum_{k\ge1}\left(\frac1k-\log(1+1/k)\right),
\]
and the alternating logarithm bounds. The tail after $K$ terms lies
between
\[
 \frac1{2(K+1)}-\frac1{6K^2}
 \quad\hbox{and}\quad\frac1{2K}.
\]
Taking $K=64$ gives a rational interval for
$H_{64}-\log65$ plus that tail. The computed enclosure lies strictly
inside $(1/2,3/5)$. Logarithms of positive rationals are enclosed by
range reduction to $1\le u\le2$ and
\[
 \log u=2\sum_{k=0}^{47}\frac{y^{2k+1}}{2k+1}+R,
 \qquad y=\frac{u-1}{u+1},\qquad
 0\le R\le\frac{2y^{97}}{97(1-y^2)}.
\]
The same formula supplies $\log2$ for the range reduction.
All arguments and finite sums are rational.

Truncate each series in \eqref{gs:eq:series} at $N=40$. On an interval
$0<\ell\le x\le u\le1/2$, the absolute omitted tails for $f,F,P$
are bounded respectively by
\begin{equation}\label{gs:eq:tail-bounds}
 T_f=\frac{2u^{N+1}}{(N+1)(1-u)},\qquad
 T_F=\frac{2u^N}{1-u},\qquad
 T_P=\frac{2Nu^{N-1}}{(1-u)^2}.
\end{equation}
These follow directly from $\zeta(n)<2$ and geometric sums.
The positive reflected series receive one-sided positive tails;
the alternating series receive symmetric tails. Their polynomial
parts are evaluated by interval Horner arithmetic. Every denominator
appearing in \eqref{gs:eq:H} is checked to have positive lower bound.

Partition $[1/16,1/2]$ into the following $256$ closed intervals:
\[
 \left[\frac1{16}+\frac{7j}{4096},
       \frac1{16}+\frac{7(j+1)}{4096}\right],
 \qquad 0\le j<256.
\]
For each entire interval, the program evaluates
\eqref{gs:eq:H} with the rigorous enclosures above. The smallest lower
endpoint obtained is exactly
\[
 \frac{2237344070433247164830307425599}
      {1267650600228229401496703205376}
 >\frac74.
\]
It occurs on $[2041/4096,1/2]$. The program tests the inequality
with integer arithmetic; the saved JSON file records every interval's
lower numerator for inspection. Thus the finite certificate proves
the asserted inequality on the complete compact interval.
\end{proof}

The endpoint proof, Lemma~\ref{gs:lem:certificate}, and reflection now
prove \eqref{gs:eq:global-derivative}. The expansions
\[
 f(x)\sim\log(1/x),\quad F(x)\sim1/x,\quad
 f(1-x)\sim\gamma x,\quad F(1-x)\sim\gamma
\]
give $J(x)\sim1/\log(1/x)$. Reciprocal symmetry gives the opposite
endpoint limit. Positivity of $J'$ proves strict monotonicity;
the endpoint limits prove surjectivity. Real analyticity and
$J'>0$ prove that its inverse is also real analytic.
This completes Theorem~\ref{gs:thm:monotonicity}.

\subsection{The proportional all-orders moment theorem}

For real $n,m>0$, define
\[
 M_{n,m}=\int_0^1 f(x)^n f(1-x)^m\,dx.
\]
These integrals converge. In particular, the following result applies
to the integer moment family in the manuscript.

\begin{theorem}[Unique proportional saddle and uniform expansion]
\label{gs:thm:proportional}
For $c>0$, put
\[
 \Phi_c(x)=c\log f(x)+\log f(1-x),\qquad
 s(c)=J^{-1}(1/c),\qquad
 A_c=-\Phi_c''(s(c)).
\]
Then $s(c)$ is the unique global maximizer of $\Phi_c$, and
\begin{equation}\label{gs:eq:curvature}
 A_c=\frac{F(1-s(c))}{f(1-s(c))}\mathcal H(s(c))>0.
\end{equation}
For every compact $C\subset(0,\infty)$ and every fixed $N\ge0$,
uniformly for $c=n/m\in C$ as $m\to\infty$,
\begin{equation}\label{gs:eq:laplace}
 M_{n,m}=f(s(c))^n f(1-s(c))^m
 \sqrt{\frac{2\pi}{mA_c}}
 \left\{\sum_{j=0}^N\frac{L_j(c)}{m^j}
                 +O_{C,N}(m^{-N-1})\right\}.
\end{equation}
The coefficient functions are real analytic on $(0,\infty)$, with
$L_0=1$ and
\begin{equation}\label{gs:eq:first-correction}
 L_1(c)=\frac{\Phi_c^{(4)}(s(c))}{8A_c^2}
       +\frac{5\Phi_c^{(3)}(s(c))^2}{24A_c^3}.
\end{equation}
Every higher coefficient is given by the finite Gaussian prescription
below.
\end{theorem}

\begin{proof}
Let $g=F/f$. The derivative has the decisive factorization
\[
 \Phi_c'(x)=g(1-x)(1-cJ(x)).
\]
Theorem~\ref{gs:thm:monotonicity} gives exactly one zero, with
positive derivative sign before the zero and negative sign after it.
Differentiating at that zero gives
\[
 \Phi_c''(s)=-c g(1-s)J'(s)
           =-g(1-s)\mathcal H(s),
\]
which proves \eqref{gs:eq:curvature} and nondegeneracy.

For completeness, the uniform endpoint control needed by Laplace's
method follows from elementary bounds. At $0<x\le1/2$,
$f(x)\le-\log x$, and \eqref{gs:eq:series} gives
\[
 f(1-x)\le\gamma x+\frac{\zeta(2)x^2}{2(1-x)}<2x.
\]
Thus as $x\downarrow0$,
\[
 \Phi_c(x)\le c\log\log(1/x)+\log(2x)\longrightarrow-\infty
\]
uniformly for $c$ in any fixed compact positive interval. Reflection
gives the same statement at $x\uparrow1$, using the positive lower
bound on $c$. Consequently all saddles stay inside a fixed compact
subinterval. By uniqueness and compactness, outside any sufficiently
small fixed neighborhood of the saddle there is a uniform positive
gap below $\Phi_c(s(c))$. Integration of that bound over an interval
of length one makes the complementary contribution exponentially
small relative to the saddle term.

On a common neighborhood of the saddle, every fixed phase derivative
is bounded uniformly in $c$, and $A_c$ is uniformly bounded below by a
positive number. Taylor-expand the analytic phase and set
$x=s(c)+u/\sqrt m$. A uniform local Gaussian bound follows from the
strictly negative second derivative and compactness. Taylor expansion
to any fixed order, with its uniformly bounded remainder, followed by
Gaussian integration proves \eqref{gs:eq:laplace}.
Explicitly, $L_j(c)$ is the coefficient of $t^{2j}$ in the normalized
Gaussian expectation, for $Z\sim N(0,A_c^{-1})$, of
\[
 \exp\left\{\sum_{k\ge3}
      \frac{\Phi_c^{(k)}(s(c))}{k!}Z^k t^{k-2}\right\}.
\]
Only finitely many derivatives occur at every order. The $t^2$
coefficient uses the fourth Gaussian moment and half the square of
the cubic term, giving \eqref{gs:eq:first-correction}. The same
finite formula proves real analyticity of the coefficient functions.
\end{proof}

\begin{corollary}[Balanced moments]\label{gs:cor:balanced}
Put
\[
 f_0=\frac12\log\pi,\qquad F_0=\gamma+2\log2,\qquad
 A_{\rm bal}=\frac{2F_0^2}{f_0^2}-\frac{\pi^2}{f_0}>0.
\]
Then the unique saddle of $M_{m,m}$ is $1/2$, and
\[
 M_{m,m}\sim
 f_0^{2m}\sqrt{\frac{2\pi}{mA_{\rm bal}}}.
\]
It has the full expansion \eqref{gs:eq:laplace}; in its first correction
the cubic phase derivative vanishes by reflection symmetry.
\end{corollary}
\begin{proof}
Use $J(1/2)=1$, $\psi(1/2)=-\gamma-2\log2$, and
$\psi_1(1/2)=\pi^2/2$ in \eqref{gs:eq:curvature}.
\end{proof}

\begin{corollary}[Saddle motion and a strictly convex limiting rate]
\label{gs:cor:envelope}
Define $\mathcal V(c)=\Phi_c(s(c))$. Then
\[
 s(1/c)=1-s(c),\qquad
 s'(c)=-\frac{g(s(c))}{A_c}<0,
\]
and
\[
 \mathcal V'(c)=\log f(s(c)),\qquad
 \mathcal V''(c)=\frac{g(s(c))^2}{A_c}>0,
 \qquad \mathcal V(c)=c\mathcal V(1/c).
\]
In particular $m^{-1}\log M_{cm,m}\to\mathcal V(c)$ uniformly on
compact positive $c$-intervals.
\end{corollary}
\begin{proof}
Reciprocal symmetry of $J$ gives the first identity. Differentiate
the stationary equation $\Phi_c'(s(c))=0$ for $s'$, and differentiate
the phase value for the two envelope identities. The reciprocity of
$\mathcal V$ follows by substituting $1-s(c)$ in $\Phi_{1/c}$.
The last assertion follows from \eqref{gs:eq:laplace}.
\end{proof}

\paragraph{What remains open after this theorem.}
The proportional regime on compact positive exponent ratios is now
resolved, including uniqueness, curvature and all fixed asymptotic
orders. A proof of the global inequality \eqref{gs:eq:global-derivative}
using only an analytic inequality, without the finite interval
certificate, remains a worthwhile simplification problem. The uniform
bridge joining this compact-$c$ regime to the logarithmic transition
of Theorem~\ref{joint:thm:transition}, and extending beyond that
window, remains a separate problem. Compact-$c$ uniformity alone
does not justify substitution of a growing exponent ratio.
